\section*{3 All You Need to Know About PPSZ: Proof of Theorem6}
We will explain the connection between PPSZ and width-bounded resolution lower bounds. After this section, the reader can forget everything about PPSZ and think of this paper as proving a certain resolution width lower bound. If $C=(C'\vee x)$ and $D=(D'\vee\bar x)$ are clauses, then $(C'\vee D')$ is called the resolvent of $C$ and $D$. It is clear that $C\wedge D$ logically implies $C'\vee D'$. Let $\Phi$ be a CNF formula. A resolution derivation from $\Phi$ is a sequence of clauses $C_1,\ldots,C_t$ such that every $C_i$ is (1) a clause of $\Phi$ or (2) the resolvent of two earlier clauses. The width of the derivation is $\max_{1\le i\le t}|C_i|$. For a clause $C$, we denote by $\operatorname{width}(\Phi\vdash C)$ the minimum width of a resolution derivation from $\Phi$ that contains $C$. Resolution is complete for refutations, that is, $\Phi$ is unsatisfiable if and only if there is a derivation of the empty clause, denoted by $\square$, from $\Phi$.
\begin{proof}[Proof of Theorem6]
Let $G,H,\phi$ be as in Theorem6, and let $\Phi:=\operatorname{TseitinBridge}(G,H,\phi)$. The only satisfying assignment of $\Phi$ is $\mathbf0$. Consider a variant of PPSZ run on $\Phi$ such that whenever it has to pick a random value for a variable, it correctly sets it to0. Fix a permutation $\pi$. Let $Y(\pi)$ be the set of variables for which this variant of PPSZ under $\pi$ could not derive the value using $R_w$, and let $Z(\pi):=\operatorname{var}(\Phi)\setminus Y(\pi)$ be the rest, i.e., all variables whose value can be derived using $R_w$ once all variables before them in $\pi$ are set to0. It is not difficult to see that the success probability of the actual PPSZ on $\Phi$ is exactly $\mathbb E_\pi[2^{-|Y(\pi)|}]$.
Suppose, for the sake of contradiction, that PPSZ using heuristic $R_w$ has success probability greater than $2^{-m+(1+\epsilon)n}$. Then there is some $\pi$ for which $Z(\pi)\ge(1+\epsilon)n$. Fix this $\pi$ and set $Z:=Z(\pi)$ and $Y:=Y(\pi)$. The set of variables $Z$ corresponds to a set $F$ of edges, $F=\{e\in E(G):x_e\in Z\}$. Set $G'=(V,F)$. Note that $G'$ has $n$ vertices and at least $(1+\epsilon)n$ edges. Setting a variable in $\Phi$ to0 corresponds to simply deleting the corresponding edge in $G$, and therefore
\[
\Phi|_{Y\mapsto0}=\operatorname{TseitinBridge}(G',H,\phi).
\]
For a graph $G=(V,E)$ and a set $X\subseteq V$, define the edge boundary $\partial(X):=\{e\in E:|e\cap X|=1\}$. Call $G$ an $(a,b)$-expander if $|\partial(X)|\ge b$ for all sets $X$ of exactly $a$ vertices. The next lemma is basically Lemma17 from [9], adapted for our purposes. We give a proof for completeness.
\begin{authorclaim}{Lemma7}
Let $\epsilon>0$ and let $G'$ be a graph on $n$ vertices with at least $(1+\epsilon)n$ edges. Let $\ell\in\mathbb N$ and $h=\ell/\epsilon$. If $h<\operatorname{girth}(G')$ then $G'$ contains a non-empty subgraph $G''$ that has minimum degree at least2 and is an $(h,\ell+1)$-expander.
\end{authorclaim}
\begin{proof}
Start with $G''=G'$. If $G''$ has a vertex of degree0 or1, delete it. If $G''$ contains a set $X$ of $h$ vertices with $|\partial(X)|\le\ell$, delete $X$ from $G''$, along with all incident edges.
The first type of deletion removes one vertex and at most one edge. The second type removes exactly $h$ vertices. There are at most $\ell$ edges in the boundary of $X$; since $|X|<\operatorname{girth}(G')$, the graph $G''[X]$ is a forest, and thus there are at most $h-1$ edges within $X$. Thus, removing $X$ removes at most $\ell+h-1<(1+\epsilon)h$ edges.
We see that a step that removes $a$ vertices removes fewer than $(1+\epsilon)a$ edges. Suppose the process terminates with $t$ vertices deleted. Trivially $t\le n$. Fewer than $(1+\epsilon)n$ edges have been deleted, so $G''$ is non-empty.
\end{proof}
Let $G''$ be given by Lemma7 with $\ell:=w+1$. We will further restrict $\Phi$ so that only edges of $G''$ remain unset. Let $F'':=E(G)\setminus E(G'')$, $Y'':=\{x_e:e\in F''\}$, and $\Phi'':=\Phi|_{Y''\mapsto0}$. Note that $\Phi''=\operatorname{TseitinBridge}(G'',H,\phi)$. Recall that all edges of $G''$ are mentioned in $Z$ and since $Y''\supseteq Y$ and restricting additional variables cannot increase the resolution width, we conclude that there exists $e\in E(G'')$ such that
\begin{equation}
\operatorname{width}(\Phi''\vdash\bar x_e)\le w.\tag{3}
\end{equation}
Towards a contradiction, we claim that in fact this resolution width is large for all variables $x_e$ where $e\in E(G'')$. Indeed, we have the following theorem:
\begin{authorclaim}{Theorem8 (Resolution Lower Bound)}
Let $G$ be a graph of minimum degree2 that is an $(h,\ell+1)$-expander. Suppose there is a locally injective homomorphism $\phi:G\to H$ into some graph $H$. Then
\begin{equation}
\operatorname{width}(\operatorname{TseitinBridge}(G,H,\phi)\vdash\bar x_e)>\ell-1,\tag{4}
\end{equation}
for all edges $e$ of $G$, provided that $2h\ell+5h+\ell<\operatorname{girth}(H)$.
\end{authorclaim}
Note that $G''$ has minimum degree2 and is a $(h,\ell+1)$-expander for $\ell=w+1$ and $h=(w+1)/\epsilon$. Also note that $\phi:V(G'')\to V(H)$ (or rather, the restriction of $\phi$ to $V(G'')$) is still a locally injective homomorphism. Recall that $w=\sqrt{\epsilon\cdot\operatorname{girth}(H)/2}-3$ and hence $2h\ell+5h+\ell=2(w+1)^2/\epsilon+5(w+1)/\epsilon+w+1<\operatorname{girth}(H)$, and thus Theorem8 applies to $G''$. This contradicts(3) and finishes the proof of Theorem6.
\end{proof}
